\documentclass{article}
\usepackage[a4paper, total={6in,8in}]{geometry}
\usepackage{graphicx,float,hyperref}
\usepackage{outlines,enumerate}
\usepackage{amsmath,amssymb,amsthm,stmaryrd}

\renewcommand{\vec}[1]{%
	\mathbf{#1}%
}

\hypersetup{
	colorlinks=true,
	linkcolor=blue,
	urlcolor=blue
}

\title{NN Exercise 1: Logistic Regression}
\date{September 25, 2026}
\author{Grey Cole}

\newcommand{\gv}{%
	\:|\:
}

\begin{document}
	\maketitle
	
	\section{Equivalence}
	Using $\exp(\mu) \cdot \exp(-\mu) = 1$,
	\begin{align*}
		\frac{\exp(\mu)}{1+\exp(\mu)} &= \frac{\exp(\mu)}{1+\exp(\mu)} \cdot \frac{\exp(-\mu)}{\exp(-\mu)}\\
		&= \frac{\exp(\mu) \cdot \exp(-\mu)}{\exp(-\mu) + \exp(\mu) \cdot \exp(-\mu)}\\
		&= \frac{1}{\exp(-\mu) + 1}
	\end{align*}
	
	In other words, dividing each part of the fraction $\frac{\exp(\mu)}{1 + \exp(\mu)}$ by $\exp(\mu)$ yields $\frac{1}{\exp(-\mu) + 1}$.
	
	\section{Inverse of Bernoulli Log Odds}
	First, we observe that, since $\theta$ is the parameter of a Bernoulli distribution, $P(0\gv  \theta) = 1 - P(1\gv \theta)$. Thus we can rewrite $f(\theta) = \log(\frac{P(1\gv \theta)}{1-P(1\gv \theta)})$.\\
	By the definition of a Bernoulli distribution with parameter $\theta$, $P(1\gv \theta) = \theta$, and with $\mu = f(\theta)$ we have 
	\begin{equation*}
		\mu = \log(\frac{P(1\gv \theta)}{1 - P(1\gv \theta)}) = \log(\frac{\theta}{1-\theta})
	\end{equation*}
	We take the exp on both sides
	\begin{equation*}
		\exp(\mu) = \frac{\theta}{1-\theta}
	\end{equation*}
	and by rearrangement we will arrive at
	\begin{align*}
		\theta &= (1-\theta)\exp(\mu)\\
		&= \exp(\mu) - \exp(\mu)\theta
	\end{align*}
	\begin{gather*}
		\theta + \exp(\mu)\theta = \exp(\mu)\\
		\theta\cdot(1 + \exp(\mu)) = \exp(\mu)\\
		\boxed{\theta = \frac{\exp(\mu)}{1 + \exp(\mu)}}
	\end{gather*}
	From Q1, we know $\frac{\exp(\mu)}{1+\exp(\mu)} = \frac{1}{\exp(-\mu) + 1}$. Thus, we have shown
	\begin{align*}
		\boxed{\theta = f^{-1}(\mu) = \frac{1}{1 + \exp(-\mu)}}
	\end{align*}
	
	\section{Total Probability}
	Since $\theta$ is a Bernoulli parameter, $P(1\gv \theta) + P(0\gv \theta) = 1$.
	\\Rewriting $P(0\gv \theta) = 1 - P(1\gv \theta)$, we have
	\begin{equation*}
		P(0\gv \theta) = 1 - \frac{1}{1+\exp(-\mu)}
	\end{equation*}
	Writing terms over a common denominator,
	\begin{equation*}
		\boxed{
			P(0\gv \theta) = \frac{1+\exp(-\mu)}{1+\exp(-\mu)} - \frac{1}{1+\exp(-\mu)}
			= \frac{\exp(-\mu)}{1+\exp(-\mu)}
		}
	\end{equation*}
	
	\section{Decision Boundary}
	If we define $-\mu = \vec{w}\vec{x}_i$, then we can use previously proven equivalencies:
	\begin{align*}
		&P(Y_i = 1 \gv  \vec{w},\vec{x}_i) = \frac{1}{1+\exp(-\mu)} = P(1\gv \theta)\\
		&P(Y_i = 0\gv  \vec{w},\vec{x}_i) = \frac{\exp(-\mu)}{1+\exp(-\mu)} = P(0\gv \theta)
	\end{align*}
	Thus we will classify an example as positive $(Y_i = 1)$ if and only if
	\begin{equation*}
		1 < \frac{P(1\gv \theta)}{P(0\gv \theta)}
	\end{equation*}
	which we further rewrite as
	\begin{equation*}
		1 < \frac{\frac{1}{1+\exp(-\mu)}}{\frac{\exp(-\mu)}{1+\exp(-\mu)}} \qquad \overset{\text{simplify}}\longrightarrow \qquad 1 < \frac{1}{\exp(-\mu)}
	\end{equation*}
	Remembering our definition $-\mu = \vec{w}\vec{x}_i$, we have a positive classification when
	\begin{equation*}
		1 < \frac{1}{\exp(\vec{w}\vec{x}_i)}
		\qquad \overset{\cdot\exp(\vec{w}\vec{x}_i)}\longrightarrow \qquad
		\exp(\vec{w}\vec{x}_i) < 1
	\end{equation*}
	Thus we have shown this is equivalent to classifying an example as positive $(Y_i = 1)$ if and only if
	\begin{equation*}
		1 > \exp(\vec{w}\vec{x}_i)
	\end{equation*}
	
	\section{Log Likelihood}
	To start, we recognize that $\vec{y}$ consists of binary target variables; that is, $y_i \in \{0,1\}$. We can rewrite the likelihood function in terms of $P(Y=1)$ and $P(Y=0)$ by utilizing $y_i$ as a selector within the product:
	\begin{equation*}
		\mathcal{L}(\vec{y}\gv x,w) = \prod_{i=1}^n P(Y_i=1\gv \vec{w},\vec{x}_i)^{y_i}\cdot P(Y_i = 0\gv \vec{w},\vec{x}_i)^{(1-y_i)}
	\end{equation*}
	Since $y_i\in{0,1}$, when $y_i=1$ the first factor is selected and when $y_i=0$ the second factor is selected, since $a^1b^0=a$ and $a^0b^1=b$.\\
	Here we will define $\theta_i = P(Y_i = 1 \gv  \vec{w},\vec{x}_i)$ as a shorthand, to make the following steps cleaner (notice that $1-\theta_i = P(Y_i = 0\gv \vec{w},\vec{x}_i)$ since $\theta_i$ is binary).
	\begin{align*}
		\ln\mathcal{L}(\vec{y}\gv x,w) &= \ln\left[ \prod_{i=1}^n \theta_i^{y_i}(1-\theta_i)^{(1-y_i)} \right]\\
		&= \sum_{i=1}^n \ln(\theta_i^{y_i}(1-\theta_i)^{(1-y_i)})\\
		&= \sum_{i=1}^n y_i\ln\theta_i + (1-y_i)\ln(1-\theta_i)\\
		&= \boxed{\sum_{i=1}^n y_i\ln P(Y_i = 1\gv w,x_i) + (1-y_i)\ln P(Y_i = 0\gv w,x_i )}\\
		& \qquad\text{Since }y_i^* \in \{0,1\},y_i = \llbracket y_i^* = 1\rrbracket \text{ and } 1 - y_i = \llbracket y_i^* = 0 \rrbracket:\\
		&= \boxed{\sum_{i=1}^n \llbracket y_i^* = 1\rrbracket\ln P(Y_i = 1\gv w,x_i) + \llbracket y_i^* = 0 \rrbracket\ln P(Y_i = 0\gv w,x_i)}
	\end{align*}
	
	
	\section{Identities}
	\subsection*{a)}
	Replace $p_0$ and $p_1$ in terms of $\vec{w}\vec{x}_i$:
	\begin{align*}
		\vec{w}\vec{x}_i &=
		\ln\left( \frac{\frac{\exp(\vec{w}\vec{x}_i)}{1+\exp(\vec{w}\vec{x}_i)}}{\frac{1}{1+\exp(\vec{w}\vec{x}_i)}} \right)\\
		&= \ln\left( \frac{\exp(\vec{w}\vec{x}_i)}{1} \right)\\
		&= \vec{w}\vec{x}_i
	\end{align*}
	Since $\vec{w}\vec{x}_i = \vec{w}\vec{x}_i$ we have shown the expression $\ln\frac{p_0}{p_1} = \vec{w}\vec{x}_i$ holds.
	
	\subsection*{b)}
	Replace $p_1$ in terms of $\vec{w}\vec{x}_i$:
	\[ \ln\left(\frac{1}{1+\exp(\vec{w}\vec{x}_i)}\right) = -\ln(1+\exp(\vec{w}\vec{x}_i)) \]
	Using logarithm rules:
	\begin{gather*} 
		\ln\left(\frac{1}{1+\exp(\vec{w}\vec{x}_i)}\right) = \ln(1) - \ln(1+\exp(\vec{w}\vec{x}_i)) \\
		\text{Then substitute:}\\
		0 - \ln(1+\exp(\vec{w}\vec{x}_i))= -\ln(1+\exp(\vec{w}\vec{x}_i))
	\end{gather*}
	Thus we have shown $\ln p_1 = \ln(1+\exp(\vec{w}\vec{x}_i))$ holds.
	
	\section{Simplified Log-Likelihood}
	First, we restate the log-likelihood function in terms of $p_0$ and $p_1$:
	\[ \ln\mathcal{L} = \sum_{i=1}^n y_i \ln p_{1i} + (1-y_i)\ln p_{0i} \]
	We can rearrange the sum:
	\begin{align*}
		y\ln p_1 + (1 - y)\ln p_0 &= y\ln p_1 + \ln p_0 - y \ln p_0\\
		&= y(\ln p_1 - \ln p_0) + \ln p_0\\
		&= y\ln\left(\frac{p_1}{p_0}\right) + \ln p_0
	\end{align*}
	Using previous identities and equivalencies:
	\begin{align*}
		y\ln\left(\frac{p_1}{p_0}\right) + \ln p_0 &=
		y \ln \left(\frac{1}{\exp(\vec{w}\vec{x})}\right) + \ln\left(\frac{\exp(\vec{w}\vec{x})}{1+\exp(\vec{w}\vec{x})}\right)\\
		&= y\cdot -\vec{w}\vec{x} + \vec{w}\vec{x} - \ln(1+\exp(\vec{w}\vec{x}))\\
		&= \boxed{(1-y)\vec{w}\vec{x} - \ln(1+\exp(\vec{w}\vec{x}))}
	\end{align*}
	By putting this expression back into the context of the log-likelihood function, we see that the log-likelihood expression does in fact simplify as:
	\[ \ln\mathcal{L}(\vec{y}\gv x,w) = \sum_{i=1}^n
		(1-y_i)\vec{w}\vec{x}_i - \ln(1 + \exp(\vec{w}\vec{x}_i)) \]
	\section{Gradient on Log-Likelihood}
	To find $\nabla(-\ln\mathcal{L}(\vec{y}\gv \vec{X},\vec{w}))$ we will differentiate with respect to $\vec{w}$ because that is the variable we are optimizing for; $\vec{y}$ and $\vec{X}$ are data points, and so can be treated as constants.\\
	First, rewrite the expression from task 7 with the negative sign:
	\begin{align*}
		-\ln\mathcal{L} &= -\sum_{i=1}^n \left[ (1-y_i)\vec{w}\vec{x}_i - \ln(1+\exp(\vec{w}\vec{x}_i)) \right]\\
		&= \sum_{i=1}^n -(1-y_i)\vec{w}\vec{x}_i + \ln(1+\exp(\vec{w}\vec{x}_i))
	\end{align*}
	Now we differentiate with respect to $w$.\\
	The first term is easy:
	\[ \nabla_{\vec{w}}(-(1-y_i)\vec{w}\vec{x}_i) = -(1-y_i)\vec{x}_i \]
	The second term requires the chain rule:
	\begin{align*}
		\nabla_{\vec{w}} (\ln1+\exp(\vec{w}\vec{x}_i)) &= \frac{1}{1+\exp(\vec{w}\vec{x}_i)} + \nabla_{\vec{w}} \exp(\vec{w}\vec{x}_i)\\
		&= \frac{1}{1+\exp(\vec{w}\vec{x}_i)} + \exp(\vec{w}\vec{x}_i)\vec{x}_i\\
		&= \frac{\exp(\vec{w}\vec{x}_i)}{1+\exp(\vec{w}\vec{x}_i)}\vec{x}_i\\
		&= p_0\vec{x}_i \qquad\qquad \leftarrow \text{Using task 6 definition}
	\end{align*}
	Putting everything back into the sum, we have the derivation of the gradient:
	\[
		\nabla_{\vec{w}} (-\ln\mathcal{L}) = \sum_{i=1}^n -(1-y_i)\vec{x}_i + p_{0i}\vec{x}_i
	\]
	Factor out $\vec{x}$, then move the negative outside the sum:
	\[
		= \sum_{i=1}^n \vec{x}_i(p_{0i} - 1 + y_i) 
		 = \boxed{- \sum_{i=1}^n \vec{x}_i(1 - y_i - p_{0i})}
	\]
	Since $(1 - p_0) = p_1$,
	\[ = -\sum_{i=1}^n \vec{x}_i(- y_i + p_{1i}) = \boxed{\sum_{i=1}^n \vec{x}_i(y_i - p_{1i})}\]
	Using definitions of $p_0$ and $p_1$, the gradient for entry $w_j$ thus has the form:
	\begin{align*}
		\frac{\partial(-\ln\mathcal{L}(\vec{y}\gv\vec{X},\vec{w}))}{\partial w_j} &= \boxed{- \sum_{i=1}^n x_{ij}((1 - y_i) - P(Y_i = 0\gv\vec{w},\vec{x}_i))}\\
		&= \boxed{\sum_{i=1}^n x_{ij}(y_i - P(Y_i = 1\gv\vec{w},\vec{x}_i)}
	\end{align*}
	
	
	\section{LogReg Gradient as Matrix Operation}
	First, we can explicitly extend the definition of $g_i(\vec{w})$ for each example $i : 1..n$ to a matrix $\vec{g}(\vec{w})$, where each row $i$ contains $g_i(\vec{w})$. Then we can work backwards from the vector matrix product
	\[ \nabla(-\ln\mathcal{L}(\vec{y}\gv\vec{X},\vec{w})) = \vec{g}(\vec{w})^T \vec{X} \] 
	The transpose $\vec{g}(\vec{w})^T$ contains 1 row and $n$ columns:
	\[ \begin{bmatrix} g_1(\vec{w}), g_2(\vec{w}), ..., g_n(\vec{w}) \end{bmatrix} \]
	The design matrix $\vec{X}$ contains $n$ rows and $1$ column (of vectors $\vec{x}$):
	\[\begin{bmatrix}
		\vec{x}_1\\
		\vec{x}_2\\
		\vdots \\
		\vec{x}_n
	\end{bmatrix}\]
	Applying matrix multiplication between $\vec{g}(\vec{w})^T$ and $\vec{X}$ yields a sum
	\begin{align*}
		 & \vec{x}_1g_1(\vec{w}) + \vec{x}_2g_2(\vec{w}) + ... + \vec{x}_ng_n(\vec{w})\\
		=& \sum_{i=1}^n x_ig_i(\vec{w})
	\end{align*}
	We can expand $g_i(\vec{w})$ back to its definition:
	\[ = \sum_{i=1}^n \vec{x}_i(y_i - P(Y_i = 1\gv \vec{w},\vec{x}_i)) \]
	However, in task 8, we have shown that this expression is equivalent to the gradient
	\[ \nabla(-\ln\mathcal{L}(\vec{y}\gv\vec{X},\vec{w}))\]
	Thus, we have shown that the gradient can be represented by a vector matrix product as
	\[ \nabla(-\ln\mathcal{L}(\vec{y}\gv\vec{X},\vec{w})) = \vec{g}(\vec{w})^T \vec{X} \]
	
	\section{L2 regularization}
	\subsection*{a)}
	As we gradually increase $\lambda$ from 0 to a large value, the regression will begin prioritizing keeping any/all $w_j$ small over keeping $-\ln\mathcal{L}$ small. Since $\|\vec{w}\|$ squares $w_j$, increasing $w_j$ will need to provide a lot more benefit the higher it is, leading to values of $w_j$ that are more similar. But for very large $\lambda$s, all $w_j$ will approach 0 as $\lambda$ approaches infinity, which diminishes any effectiveness altogether.
	\subsection*{b)}
	\begin{align*}
		\frac{\partial}{\partial w_j}\mathcal{O}(\vec{w}) &= \frac{\partial}{\partial w_j}(-\ln\mathcal{L}(\vec{y}\gv\vec{X},\vec{w})) + \frac{\partial}{\partial w_j}(\frac{\lambda}{2}\sum_j w_j^2)\\
		&= \boxed{ \sum_{i=1}^n x_{ij}g_i(\vec{w}) + \lambda w_j }
	\end{align*}
	The partial derivative $\frac{\partial}{\partial w_j}\mathcal{O}(\vec{w})$ applies a flat bias $\lambda w_j$ to each dimension in the weight vector, while $\mathcal{O}(\vec{w})$ applies the sum of each weight squared to the regularized NLL. L2 Regularization is likely referred to as weight decay because it pushes weights toward 0, by imposing a penalty on larger weights, especially across the entire NLL from squared weights.
	
	%\section{$\star$ L2 regularization as Gaussian Prior}
	
	%\section{$\star$ Separable data}
	
	%\section{$\star$ Convexity}
\end{document}